\chapter{Implicaciones sociales y éticas del proyecto}
\label{ch:anexo_etica}

La seguridad no es un concepto abstracto. Detrás de cada dispositivo conectado
a una red hay una persona, una familia, una institución o una comunidad que
confía en que sus datos, su identidad y su integridad estarán protegidos. Cuando
un ingeniero diseña un módulo de seguridad para sistemas embebidos, no está
manipulando compuertas lógicas y registros: está asumiendo una responsabilidad
directa sobre la confianza que otros depositarán en ese hardware. Este proyecto
me confrontó con esa realidad desde el primer día.

El Internet de las Cosas ha transformado la vida cotidiana de formas que hace
dos décadas habrían parecido ciencia ficción. Dispositivos médicos que monitorean
signos vitales, sensores industriales que controlan procesos críticos, cerraduras
inteligentes que protegen hogares y sistemas vehiculares que toman decisiones en
fracciones de segundo. Todos ellos comparten una característica común: operan en
entornos físicamente accesibles, donde un atacante con las herramientas adecuadas
puede intentar extraer claves, clonar identidades o manipular el comportamiento
del dispositivo. La consecuencia de una falla de seguridad no se mide solo en
pérdidas económicas, sino en riesgos para la salud, la privacidad y la seguridad
de las personas.

Las Funciones Físicamente No Clonables representan un cambio de paradigma en la
forma de abordar la seguridad del hardware. Al eliminar la necesidad de almacenar
claves criptográficas en memorias que pueden ser leídas o copiadas, las PUFs
ofrecen una capa de protección fundamentada en la física misma del silicio,
imposible de replicar incluso por el fabricante original. Esta promesa tecnológica
conlleva una responsabilidad ética significativa. Quien diseña una PUF debe ser
riguroso en la evaluación de sus propiedades estadísticas, honesto sobre sus
limitaciones y consciente de que una evaluación incompleta puede generar una falsa
sensación de seguridad con consecuencias graves.

Durante el desarrollo de este proyecto, una de las decisiones más difíciles fue
aceptar y documentar abiertamente las limitaciones de la evaluación realizada. La
ausencia de pruebas inter-dispositivo impide afirmar con certeza que el módulo
PLPUF implementado sea capaz de distinguir chips individuales en un escenario de
autenticación real. Habría sido tentador omitir esta limitación o minimizarla,
pero hacerlo habría comprometido la integridad del trabajo y la seguridad de
quienes en el futuro pudieran utilizar este módulo. La honestidad intelectual no
es un lujo en ingeniería de seguridad: es un requisito ético fundamental.

Otro aspecto que merece reflexión es la decisión de desarrollar este proyecto
como código abierto. En el ámbito de la seguridad informática existe un debate
permanente entre quienes abogan por la seguridad mediante la privacidad y quienes
defienden la transparencia. Mi convicción, respaldada por décadas de experiencia,
es que un diseño de seguridad solo es confiable tras el escrutinio público. Al
publicar el código HDL, los controladores y los datos de evaluación, este proyecto
invita a la comunidad a verificar, reproducir y mejorar los resultados. Esta
apertura no debilita la seguridad; la fortalece.

No obstante, la publicación de conocimientos sobre seguridad del hardware plantea
el dilema del uso dual. El mismo conocimiento que permite diseñar una PUF robusta
puede emplearse para explotar sus debilidades. Este dilema es inherente a toda
investigación en seguridad. Creo que la contribución neta del conocimiento abierto
es positiva: los defensores necesitan más información que los atacantes, y la
transparencia permite corregir vulnerabilidades antes de que sean explotadas.

Desde una perspectiva más amplia, este proyecto contribuye a la democratización
de la seguridad del hardware. Las soluciones comerciales suelen estar protegidas
por patentes que las hacen inaccesibles para instituciones educativas, pequeñas
empresas y países en desarrollo. Al proporcionar una plataforma completa,
documentada y gratuita, se abre la posibilidad de que estudiantes e investigadores
de todo el mundo experimenten con técnicas de seguridad sin costosas licencias.
En Costa Rica, donde formar ingenieros capaces de diseñar sistemas seguros es una
necesidad estratégica, esta contribución adquiere un valor que trasciende lo
técnico.

Finalmente, deseo reflexionar sobre el rol del ingeniero en la sociedad digital
contemporánea. Vivimos en una era donde la confianza en la tecnología es un pilar
de la convivencia social. Cada vez que alguien usa un dispositivo conectado,
deposita su confianza en los ingenieros que lo diseñaron. Esta confianza no debe
tomarse a la ligera. Mi deber como ingeniero no se agota en implementar un
circuito u optimizar una métrica: se extiende a garantizar que mi trabajo
contribuya al bienestar de las personas y no se convierta en un instrumento de
vulnerabilidad. Con este proyecto, he intentado honrar esa responsabilidad con
rigor técnico, transparencia y compromiso con la comunidad.